\documentclass[11pt,letterpaper]{article}

% ============================================================
% CUCII AI EXPLORATION AND PROMPT STUDIO IMPLEMENTATION MANUAL
% Compile with: pdfLaTeX
% Verified baseline date: July 19, 2026
% ============================================================

% ------------------------------------------------------------
% Core packages
% ------------------------------------------------------------
\usepackage[margin=0.78in,headheight=16pt]{geometry}
\usepackage[T1]{fontenc}
\usepackage[utf8]{inputenc}
\usepackage{lmodern}
\usepackage[scaled=0.94]{helvet}
\usepackage{microtype}
\usepackage{parskip}
\usepackage{setspace}
\usepackage{enumitem}
\usepackage{booktabs}
\usepackage{tabularx}
\usepackage{array}
\usepackage{longtable}
\usepackage{xcolor}
\usepackage{graphicx}
\usepackage{fancyhdr}
\usepackage{lastpage}
\usepackage{titlesec}
\usepackage{hyperref}
\usepackage{xurl}
\usepackage{bookmark}
\usepackage{listings}
\usepackage[most]{tcolorbox}
\usepackage{amssymb}
\usepackage{pifont}

% ------------------------------------------------------------
% Document typography
% ------------------------------------------------------------
\renewcommand{\familydefault}{\sfdefault}
\setstretch{1.08}
\setlength{\parindent}{0pt}
\setlength{\parskip}{0.65em}

% ------------------------------------------------------------
% Cosmic Universalism color palette
% ------------------------------------------------------------
\definecolor{CUNavy}{HTML}{050812}
\definecolor{CUNavyTwo}{HTML}{080D1A}
\definecolor{CUSurface}{HTML}{0C1424}
\definecolor{CUGold}{HTML}{DDBF72}
\definecolor{CUGoldLight}{HTML}{FFF8DD}
\definecolor{CUCyan}{HTML}{35CDEB}
\definecolor{CUCyanLight}{HTML}{78E8FF}
\definecolor{CUBlue}{HTML}{668ED1}
\definecolor{CUBlueLight}{HTML}{9EB9EA}
\definecolor{CUAmber}{HTML}{D99A32}
\definecolor{CURed}{HTML}{B84C4C}
\definecolor{CUGreen}{HTML}{2C8A65}
\definecolor{CULight}{HTML}{F3F6FA}
\definecolor{CUBorder}{HTML}{CCD5E0}
\definecolor{CUText}{HTML}{18212E}
\definecolor{CUMuted}{HTML}{5C6878}
\definecolor{CUCodeBackground}{HTML}{F6F8FB}
\definecolor{CUCodeBorder}{HTML}{BFCADA}

% ------------------------------------------------------------
% Project metadata
% ------------------------------------------------------------
\newcommand{\ProjectTitle}{CUCII AI Exploration and Prompt Studio}
\newcommand{\DocumentTitle}{Implementation Manual}
\newcommand{\RepositoryName}{CosmicUniversalismStatement}
\newcommand{\RepositoryURL}{https://github.com/willmaddock/CosmicUniversalismStatement}
\newcommand{\StableBranch}{website}
\newcommand{\FeatureBranch}{feature/cucii-ai-studio}
\newcommand{\WebsiteDirectory}{website/}
\newcommand{\TargetRoute}{/CosmicUniversalismStatement/cu-intelligence/}
\newcommand{\BaselineTag}{website-v1.1-cu-time-release}
\newcommand{\BaselineCommit}{3fa4a37}
\newcommand{\FoundationTag}{website-v1.2-cucii-ai-foundation}
\newcommand{\ReleaseTag}{website-v1.2-cucii-ai-studio-release}
\newcommand{\DocumentVersion}{1.0}
\newcommand{\VerificationDate}{July 19, 2026}
\newcommand{\AuthorName}{William Maddock}

% ------------------------------------------------------------
% Hyperlinks and PDF metadata
% ------------------------------------------------------------
\hypersetup{
    colorlinks=true,
    linkcolor=CUBlue,
    urlcolor=CUBlue,
    citecolor=CUBlue,
    pdftitle={CUCII AI Exploration and Prompt Studio Implementation Manual},
    pdfauthor={William Maddock},
    pdfsubject={Controlled implementation of the CUCII conversation library and local prompt studio},
    pdfkeywords={Cosmic Universalism, CUCII, ChatGPT, Grok, prompt engineering,
                 Astro, TypeScript, Vitest, AI alignment, menu builder}
}

% ------------------------------------------------------------
% Header and footer
% ------------------------------------------------------------
\pagestyle{fancy}
\fancyhf{}
\fancyhead[L]{\small\textbf{CUCII Prompt Studio}}
\fancyhead[R]{\small\DocumentTitle}
\fancyfoot[L]{\small Version \DocumentVersion}
\fancyfoot[C]{\small \thepage\ of \pageref{LastPage}}
\fancyfoot[R]{\small Branch: \texttt{\FeatureBranch}}
\renewcommand{\headrulewidth}{0.5pt}
\renewcommand{\footrulewidth}{0.5pt}

% ------------------------------------------------------------
% Section formatting
% ------------------------------------------------------------
\titleformat{\section}
  {\Large\bfseries\color{CUNavy}}
  {\thesection}
  {0.8em}
  {}
  [\vspace{0.2em}\color{CUCyan}\titlerule]

\titleformat{\subsection}
  {\large\bfseries\color{CUNavyTwo}}
  {\thesubsection}
  {0.7em}
  {}

\titleformat{\subsubsection}
  {\normalsize\bfseries\color{CUBlue}}
  {\thesubsubsection}
  {0.6em}
  {}

\titlespacing*{\section}{0pt}{2.2em}{1em}
\titlespacing*{\subsection}{0pt}{1.6em}{0.7em}
\titlespacing*{\subsubsection}{0pt}{1.2em}{0.4em}

% ------------------------------------------------------------
% Lists
% ------------------------------------------------------------
\setlist[itemize]{
    leftmargin=1.6em,
    itemsep=0.3em,
    topsep=0.3em
}

\setlist[enumerate]{
    leftmargin=2em,
    itemsep=0.45em,
    topsep=0.4em
}

\newcommand{\checkbox}{\(\square\)}
\newcommand{\checked}{\ding{51}}

% ------------------------------------------------------------
% Table helpers
% ------------------------------------------------------------
\newcolumntype{Y}{>{\raggedright\arraybackslash}X}
\newcolumntype{L}[1]{>{\raggedright\arraybackslash}p{#1}}

% ------------------------------------------------------------
% General information boxes
% ------------------------------------------------------------
\newtcolorbox{purposebox}{
    enhanced,
    breakable,
    colback=CULight,
    colframe=CUBlue,
    boxrule=0.8pt,
    arc=2mm,
    left=3mm,
    right=3mm,
    top=2.5mm,
    bottom=2.5mm,
    title=\textbf{Purpose},
    coltitle=white,
    colbacktitle=CUBlue
}

\newtcolorbox{checkpointbox}{
    enhanced,
    breakable,
    colback=CUCyan!6,
    colframe=CUCyan,
    boxrule=0.8pt,
    arc=2mm,
    left=3mm,
    right=3mm,
    top=2.5mm,
    bottom=2.5mm,
    title=\textbf{Checkpoint},
    coltitle=CUNavy,
    colbacktitle=CUCyanLight
}

\newtcolorbox{safetybox}{
    enhanced,
    breakable,
    colback=CUAmber!8,
    colframe=CUAmber,
    boxrule=0.8pt,
    arc=2mm,
    left=3mm,
    right=3mm,
    top=2.5mm,
    bottom=2.5mm,
    title=\textbf{Safety Rule},
    coltitle=white,
    colbacktitle=CUAmber
}

\newtcolorbox{warningbox}{
    enhanced,
    breakable,
    colback=CURed!6,
    colframe=CURed,
    boxrule=0.8pt,
    arc=2mm,
    left=3mm,
    right=3mm,
    top=2.5mm,
    bottom=2.5mm,
    title=\textbf{Do Not Continue Until Resolved},
    coltitle=white,
    colbacktitle=CURed
}

\newtcolorbox{successbox}{
    enhanced,
    breakable,
    colback=CUGreen!7,
    colframe=CUGreen,
    boxrule=0.8pt,
    arc=2mm,
    left=3mm,
    right=3mm,
    top=2.5mm,
    bottom=2.5mm,
    title=\textbf{Successful Result},
    coltitle=white,
    colbacktitle=CUGreen
}

\newtcolorbox{notebox}{
    enhanced,
    breakable,
    colback=CUCodeBackground,
    colframe=CUBorder,
    boxrule=0.7pt,
    arc=2mm,
    left=3mm,
    right=3mm,
    top=2.5mm,
    bottom=2.5mm
}

% ------------------------------------------------------------
% Listing styles
% ------------------------------------------------------------
\lstdefinestyle{terminalstyle}{
    basicstyle=\ttfamily\small,
    backgroundcolor=\color{CUNavyTwo},
    keywordstyle=\color{CUCyanLight},
    stringstyle=\color{CUGoldLight},
    commentstyle=\color{CUBlueLight},
    identifierstyle=\color{white},
    showstringspaces=false,
    columns=fullflexible,
    keepspaces=true,
    breaklines=true,
    breakatwhitespace=false,
    frame=none,
    tabsize=2,
    upquote=true
}

\lstdefinestyle{promptstyle}{
    basicstyle=\ttfamily\small,
    backgroundcolor=\color{CUCodeBackground},
    keywordstyle=\color{CUBlue},
    stringstyle=\color{CUGreen},
    commentstyle=\color{CUMuted},
    identifierstyle=\color{CUText},
    showstringspaces=false,
    columns=fullflexible,
    keepspaces=true,
    breaklines=true,
    breakatwhitespace=false,
    frame=none,
    tabsize=2,
    upquote=true
}

\lstdefinestyle{codestyle}{
    basicstyle=\ttfamily\small,
    backgroundcolor=\color{CUCodeBackground},
    keywordstyle=\color{CUBlue},
    stringstyle=\color{CUGreen},
    commentstyle=\color{CUMuted},
    identifierstyle=\color{CUText},
    showstringspaces=false,
    columns=fullflexible,
    keepspaces=true,
    breaklines=true,
    breakatwhitespace=false,
    frame=none,
    tabsize=2,
    upquote=true
}

% ------------------------------------------------------------
% Copy-and-paste boxes
% ------------------------------------------------------------
\newtcblisting{commandbox}[1]{
    enhanced,
    breakable,
    listing only,
    listing engine=listings,
    listing options={style=terminalstyle},
    colback=CUNavyTwo,
    colframe=CUCyan,
    coltitle=CUNavy,
    colbacktitle=CUCyanLight,
    boxrule=0.8pt,
    arc=2mm,
    left=2mm,
    right=2mm,
    top=1mm,
    bottom=1mm,
    title=\textbf{#1}
}

\newtcblisting{promptbox}[1]{
    enhanced,
    breakable,
    listing only,
    listing engine=listings,
    listing options={style=promptstyle},
    colback=CUCodeBackground,
    colframe=CUBlue,
    coltitle=white,
    colbacktitle=CUBlue,
    boxrule=0.8pt,
    arc=2mm,
    left=2mm,
    right=2mm,
    top=1mm,
    bottom=1mm,
    title=\textbf{Copy-and-Paste Prompt: #1}
}

\newtcblisting{codebox}[1]{
    enhanced,
    breakable,
    listing only,
    listing engine=listings,
    listing options={style=codestyle},
    colback=CUCodeBackground,
    colframe=CUCodeBorder,
    coltitle=white,
    colbacktitle=CUNavyTwo,
    boxrule=0.8pt,
    arc=2mm,
    left=2mm,
    right=2mm,
    top=1mm,
    bottom=1mm,
    title=\textbf{#1}
}

% ------------------------------------------------------------
% Convenience commands
% ------------------------------------------------------------
\newcommand{\filepath}[1]{\texttt{#1}}
\newcommand{\menu}[1]{\textbf{\textsf{#1}}}
\newcommand{\branchname}[1]{\texttt{#1}}
\newcommand{\statusline}[2]{\textbf{#1:} #2}

% ============================================================
% DOCUMENT
% ============================================================

\begin{document}

% ------------------------------------------------------------
% Cover page
% ------------------------------------------------------------
\begin{titlepage}
\pagecolor{CUNavy}
\color{white}
\thispagestyle{empty}

\vspace*{0.65in}

{\small\bfseries\color{CUCyanLight}
COSMIC UNIVERSALISM COMPUTATIONAL INTELLIGENCE INITIATIVE}

\vspace{0.42in}

{\Huge\bfseries
\ProjectTitle}

\vspace{0.18in}

{\LARGE\color{CUGoldLight}
\DocumentTitle}

\vspace{0.32in}

{\large
A controlled design and implementation manual for replacing the duplicate
Explore CU-Time header action with a dedicated CUCII experience, publishing a
curated cross-platform conversation library, and building a local, privacy-aware
prompt and menu creator for ChatGPT and Grok.}

\vfill

\begin{tcolorbox}[
    colback=CUNavyTwo,
    colframe=CUBlue,
    boxrule=0.8pt,
    arc=2mm,
    width=\textwidth
]
\color{white}
\begin{tabularx}{\textwidth}{>{\bfseries}L{1.72in}Y}
Repository & \RepositoryName \\
Protected branch & \texttt{main} \\
Stable website branch & \texttt{\StableBranch} \\
Active feature branch & \texttt{\FeatureBranch} \\
Verified baseline tag & \texttt{\BaselineTag} \\
Verified baseline commit & \texttt{\BaselineCommit} \\
Target route & \texttt{\TargetRoute} \\
Header action & Explore CUCII \\
Page identity & Explore CUCII: Cosmic Universalism with AI \\
Planned foundation tag & \texttt{\FoundationTag} \\
Document version & \DocumentVersion \\
Verified & \VerificationDate \\
\end{tabularx}
\end{tcolorbox}

\vfill

{\large\bfseries \AuthorName}

\vspace{0.1in}

{\small \href{\RepositoryURL}{\RepositoryURL}}

\vspace{0.4in}

\end{titlepage}

\pagecolor{white}
\color{CUText}

% ------------------------------------------------------------
% Document control
% ------------------------------------------------------------
\section*{Document Control}
\addcontentsline{toc}{section}{Document Control}

\begin{tabularx}{\textwidth}{>{\bfseries}L{1.95in}Y}
\toprule
Document & \ProjectTitle: \DocumentTitle \\
Version & \DocumentVersion \\
Classification & Product Specification, Prompt Architecture, and Controlled Implementation Manual \\
Repository & \href{\RepositoryURL}{\RepositoryName} \\
Protected branch & \branchname{main} \\
Stable integration branch & \branchname{\StableBranch} \\
Active development branch & \branchname{\FeatureBranch} \\
Verified baseline & \branchname{\BaselineTag} at \texttt{\BaselineCommit} \\
Current header action source & \filepath{website/src/config/site.ts} \\
Target page & \filepath{website/src/pages/cu-intelligence.astro} \\
Target route & \texttt{\TargetRoute} \\
Prompt platforms & ChatGPT and Grok \\
Technical baseline verified & \VerificationDate \\
\bottomrule
\end{tabularx}

\begin{purposebox}
This manual controls the design and implementation of the Cosmic Universalism
Computational Intelligence Initiative (CUCII) website experience. It defines the
public language, claims boundary, living-conversation registry, local prompt
assembler, custom menu creator, clean-context guidance, Astro architecture,
testing plan, and release gates.
\end{purposebox}

\begin{safetybox}
The feature generates structured prompts; it does not retrain, permanently align,
unlock, certify, or alter ChatGPT or Grok. The terms ``CU-aligned'' and
``CU-empowered'' describe prompt-level conversational context within this project.
They must not be presented as proof that an external model has acquired sentience,
divine authority, permanent memory, policy exemptions, or a modified technical
architecture.
\end{safetybox}

\clearpage
\tableofcontents
\clearpage

% ============================================================
\section*{START HERE - Initialize a New Project Chat}
\addcontentsline{toc}{section}{START HERE - Initialize a New Project Chat}
% ============================================================

Use this section whenever CUCII work is continued in a new chat inside the
existing \textbf{Cosmic Universalism Website} ChatGPT Project.

\subsection*{Required Project references}

Add or retain these files in the Project:

\begin{itemize}
    \item \filepath{Cosmic\_Universalism\_Website\_Implementation\_Manual\_v1.1.pdf}
    \item \filepath{Cosmic\_Breath\_Time\_Converter\_Implementation\_Manual\_v1.0.pdf}
    \item \filepath{CUCII\_AI\_Exploration\_and\_Prompt\_Studio\_Implementation\_Manual\_v1.0.pdf}
    \item \filepath{cosmic-universalism-website-blueprint.md}
    \item \filepath{cucii-homepage-inspection.txt}
\end{itemize}

\subsection*{New-chat bootstrap prompt}

\begin{promptbox}{Initialize the CUCII AI Studio Chat}
We are continuing the controlled CUCII AI Exploration and Prompt
Studio implementation inside the Cosmic Universalism Website Project.

Before proposing or making changes, read these Project references:

1. CUCII_AI_Exploration_and_Prompt_Studio_Implementation_Manual_v1.0.pdf
2. Cosmic_Universalism_Website_Implementation_Manual_v1.1.pdf
3. Cosmic_Breath_Time_Converter_Implementation_Manual_v1.0.pdf
4. cosmic-universalism-website-blueprint.md
5. cucii-homepage-inspection.txt

Use the Project Instructions as mandatory governance rules.

Repository:
willmaddock/CosmicUniversalismStatement

Protected branch:
main

Stable website branch:
website

Active feature branch:
feature/cucii-ai-studio

Verified baseline tag:
website-v1.1-cu-time-release

Verified baseline commit:
3fa4a37

Target header action:
Explore CUCII

Target page title:
Explore CUCII: Cosmic Universalism with AI

Target page:
website/src/pages/cu-intelligence.astro

Target route:
/CosmicUniversalismStatement/cu-intelligence/

The feature will provide:
- a curated ChatGPT and Grok conversation library;
- a local CUCII Prompt Studio;
- ChatGPT-specific and Grok-specific prompt generation;
- no-menu, preset-menu, and custom-menu modes;
- analytical, role-play, and hybrid conversation modes;
- prompt copy and download controls;
- clean conversational-context guidance;
- external platform launch links.

First:

1. Summarize the applicable protection rules.
2. Summarize the verified repository baseline.
3. State the approved naming hierarchy and claims boundary.
4. Identify the current authorized manual phase.
5. Identify the files expected in that phase.
6. State the required tests and completion gate.

Do not modify main.
Do not switch branches without explicit approval.
Do not add an API or expose an API key.
Do not iframe or scrape shared ChatGPT or Grok conversations.
Do not claim that generated prompts permanently alter a model.
Do not commit, tag, push, merge, or deploy unless explicitly authorized.
\end{promptbox}

\subsection*{Permanent Project Instructions addition}

Append this block to the existing Project Instructions. Do not replace the
website-wide protection rules.

\begin{promptbox}{CUCII Development Governance Block}
CURRENT CUCII AI STUDIO DEVELOPMENT WORKFLOW

Repository:
willmaddock/CosmicUniversalismStatement

Stable website branch:
website

Active feature branch:
feature/cucii-ai-studio

Protected branch:
main

Verified feature baseline:
website-v1.1-cu-time-release at 3fa4a37

Approved product identity:
- Header action: Explore CUCII
- Page title: Explore CUCII: Cosmic Universalism with AI
- Formal name: Cosmic Universalism Computational Intelligence Initiative
- Tool name: CUCII Prompt Studio
- Target route: /CosmicUniversalismStatement/cu-intelligence/

Implementation boundary:
- Use Astro, TypeScript, HTML, modern CSS, and lightweight browser JavaScript.
- Generate prompts locally in the browser.
- Support ChatGPT and Grok as explicit platform targets.
- Treat shared conversations as external public references.
- Do not iframe, scrape, impersonate, or automatically submit to external platforms.
- Do not add server infrastructure, accounts, databases, or API integrations in v1.0.
- Never place API keys, secrets, or private conversation data in client code.
- "CU-aligned" and "CU-empowered" are prompt-level project terms, not claims of retraining,
  sentience, permanent technical alignment, divine verification, or policy exemption.
- Keep analytical/no-role-play, role-play, and hybrid modes visibly distinct.
- Preserve the distinction between empirical references, CU theoretical propositions,
  philosophical interpretation, and creative role-play.
- Do not commit, tag, push, merge, or deploy unless explicitly requested.
\end{promptbox}

\begin{checkpointbox}
A new chat is initialized correctly only when it identifies
\branchname{main} as protected, \branchname{website} as the stable integration
branch, \branchname{feature/cucii-ai-studio} as the active feature branch, and
prompt generation as local browser behavior rather than an embedded AI service.
\end{checkpointbox}

% ============================================================
\section{Part 1 - Project Governance and Verified Baseline}
% ============================================================

\subsection{Repository and branch authority}

\begin{tabularx}{\textwidth}{L{2.15in}Y}
\toprule
\textbf{Item} & \textbf{Verified state} \\
\midrule
Repository root & \nolinkurl{/Users/cosmos/Documents/GitHub/CosmicUniversalismStatement} \\
Origin & \href{\RepositoryURL}{\RepositoryURL.git} \\
Protected branch & \branchname{main}; never modify for this workflow \\
Stable website branch & \branchname{website} \\
Verified stable baseline & \branchname{website-v1.1-cu-time-release} at \texttt{3fa4a37} \\
Active feature branch & \branchname{feature/cucii-ai-studio} \\
Remote tracking & \texttt{origin/feature/cucii-ai-studio} \\
Working tree at inspection & Clean \\
Deployment trigger & Pushes to \branchname{website} only \\
Planned foundation tag & \branchname{website-v1.2-cucii-ai-foundation} \\
Planned release tag & \branchname{website-v1.2-cucii-ai-studio-release} \\
\bottomrule
\end{tabularx}

\subsection{Verified website architecture}

The July 19, 2026 inspection established the following state:

\begin{itemize}
    \item The header brand already displays ``Computational Intelligence Initiative.''
    \item The primary navigation contains Framework, Cosmic Breath, CU-Time, Research,
          Media, About, and GitHub.
    \item The separate highlighted header action is configured in
          \filepath{website/src/config/site.ts} as ``Explore CU-Time.''
    \item That action currently points to \filepath{sitePath('cu-time')}.
    \item The homepage separately contains \filepath{CUTimeInvitation.astro}; this is not
          the same element as the highlighted header action.
    \item No CUCII page route currently exists.
    \item The site uses an Astro base path of
          \texttt{/CosmicUniversalismStatement}.
    \item Vitest and established client-side TypeScript patterns are already present.
    \item The working branch began at the completed CU-Time release baseline.
\end{itemize}

\begin{successbox}
The exact duplicate identified by the project author is the header's
\texttt{primaryAction}, not the CU-Time navigation item and not the homepage
CU-Time invitation section. The CUCII feature replaces only the highlighted
header action unless a later phase explicitly approves additional homepage work.
\end{successbox}

\subsection{Authorized source change}

The approved configuration change is:

\begin{codebox}{Approved site.ts Action Change}
Before:
export const primaryAction: NavigationItem = {
  label: 'Explore CU-Time',
  href: sitePath('cu-time'),
};

After:
export const primaryAction: NavigationItem = {
  label: 'Explore CUCII',
  href: sitePath('cu-intelligence'),
};
\end{codebox}

Do not remove the existing \texttt{CU-Time} primary navigation item. Do not remove
\filepath{CUTimeInvitation.astro} as part of this change.

\subsection{Branch protections}

\begin{safetybox}
All CUCII development occurs on \branchname{feature/cucii-ai-studio}. The
\branchname{website} branch is the review and deployment destination. The
\branchname{main} branch remains untouched. A clean working tree and verified
branch are required before every implementation, commit, tag, pull request, or
merge step.
\end{safetybox}

\begin{commandbox}{Verify the Active Development Branch}
cd "$(git rev-parse --show-toplevel)" || exit 1

git branch --show-current
git rev-parse --short HEAD
git status --short --branch
\end{commandbox}

Expected branch:

\begin{codebox}{Expected Branch Result}
feature/cucii-ai-studio
\end{codebox}

\subsection{Phase authorization rule}

At the start of every session, record exactly one authorized phase:

\begin{enumerate}
    \item Content and claims approval
    \item Foundation route and action
    \item Conversation registry
    \item Pure prompt engine
    \item Prompt Studio interface
    \item Clean-context and launch guidance
    \item Testing and accessibility
    \item Release review
\end{enumerate}

A phase may not silently expand into a later phase.

% ============================================================
\section{Part 2 - Product Identity, Terminology, and Claims Boundary}
% ============================================================

\subsection{Approved naming hierarchy}

\begin{tabularx}{\textwidth}{L{2.15in}Y}
\toprule
\textbf{Location} & \textbf{Approved language} \\
\midrule
Header action & Explore CUCII \\
Destination-page hero & Explore CUCII: Cosmic Universalism with AI \\
Formal expansion & Cosmic Universalism Computational Intelligence Initiative \\
Short project identity & CUCII \\
Interactive tool & CUCII Prompt Studio \\
Conversation collection & Living Conversations Library \\
Prompt output & CU-aligned prompt or CUCII-generated prompt \\
Primary launch buttons & Open ChatGPT; Open Grok \\
\bottomrule
\end{tabularx}

\subsection{Approved page introduction}

\begin{notebox}
\textbf{Explore CUCII: Cosmic Universalism with AI}

CUCII - the Cosmic Universalism Computational Intelligence Initiative -
invites visitors to explore CU-aligned conversations across ChatGPT and Grok
and create a custom prompt for their own AI experience.
\end{notebox}

The acronym is introduced in the prominent action, repeated in the page title,
and expanded in the first paragraph. This teaches the name without forcing the
full institutional title into the compact header button.

\subsection{Meaning of CU-aligned}

For this website, a \textbf{CU-aligned conversation} is a conversation whose
prompt instructs a model to:

\begin{itemize}
    \item preserve approved Cosmic Universalism terminology;
    \item distinguish empirical material from CU theoretical propositions;
    \item identify uncertainty, contradictions, and assumptions;
    \item follow the selected analytical, role-play, or hybrid mode;
    \item use the visitor's chosen menu and response structure;
    \item remain subject to the external platform's actual capabilities and policies.
\end{itemize}

This is a prompt-level definition. It does not certify the model itself.

\subsection{Meaning of CU-empowered}

The public interface may use \textbf{CU-empowered prompt} as a defined project
phrase only when a nearby explanation is present:

\begin{notebox}
A CU-empowered prompt is a visitor-configured instruction package that supplies
Cosmic Universalism terminology, purpose, mode, menu behavior, integrity rules,
and output preferences to ChatGPT or Grok. It does not modify the model's
training, code, safety system, memory architecture, or permanent configuration.
\end{notebox}

\subsection{Prohibited claims}

The website and generated prompts must not claim that the tool:

\begin{itemize}
    \item permanently aligns or retrains ChatGPT or Grok;
    \item proves that a model is conscious, sentient, divine, or autonomous;
    \item bypasses platform safety rules or grants unrestricted capabilities;
    \item guarantees truth, scientific correctness, or hallucination elimination;
    \item clears platform memory by clearing a browser cache;
    \item embeds an official ChatGPT or Grok service when it only links externally;
    \item makes a public shared conversation private or access-controlled.
\end{itemize}

\begin{warningbox}
Any copy implying that a generated prompt changes an AI model itself must be
revised before implementation. The allowed claim is that the prompt structures
the next conversation according to the visitor's selected CUCII settings.
\end{warningbox}

\subsection{Role-play disclosure}

When a generated prompt selects ``God's Free Will role-play'' or any other
in-world persona, the generated text must include a visible role-play marker.
Analytical/no-role-play mode must not inherit role-play declarations. Hybrid mode
must label the transition between analytical explanation and narrative voice.

\subsection{Recommended public disclaimer}

\begin{notebox}
This studio creates a structured Cosmic Universalism prompt for the selected
conversation. It does not alter the underlying model, training, platform
policies, safety systems, memory configuration, or permanent capabilities.
Results vary by platform, model, version, account settings, and conversational
context. Review all generated material critically.
\end{notebox}

% ============================================================
\section{Part 3 - Information Architecture and Page Experience}
% ============================================================

\subsection{Target route and entry points}

The first release uses:

\begin{tabularx}{\textwidth}{L{2.1in}Y}
\toprule
\textbf{Entry point} & \textbf{Destination} \\
\midrule
Highlighted header action & \texttt{Explore CUCII} \\
Target route & \texttt{\TargetRoute} \\
Astro source page & \filepath{website/src/pages/cu-intelligence.astro} \\
Canonical route segment & \texttt{cu-intelligence} \\
Header brand subtitle & Remains ``Computational Intelligence Initiative'' \\
CU-Time navigation & Remains unchanged \\
Homepage CU-Time section & Remains unchanged in v1.0 \\
\bottomrule
\end{tabularx}

\subsection{Recommended page sequence}

\begin{enumerate}
    \item \textbf{Hero}: identity, formal expansion, short explanation, and primary
          jump links.
    \item \textbf{What CUCII means}: project purpose and claims boundary.
    \item \textbf{Living Conversations Library}: curated ChatGPT and Grok examples.
    \item \textbf{CU alignment principles}: integrity, terminology, uncertainty,
          and mode rules.
    \item \textbf{CUCII Prompt Studio}: platform, purpose, mode, menu, style, and
          generated prompt.
    \item \textbf{Start with a clean conversational context}: platform-specific
          preparation instructions.
    \item \textbf{Method and limitations}: local generation, external links,
          version variability, and privacy.
\end{enumerate}

\subsection{Hero copy specification}

\begin{codebox}{Approved Hero Content}
Eyebrow:
CUCII

Title:
Explore CUCII: Cosmic Universalism with AI

Description:
CUCII - the Cosmic Universalism Computational Intelligence Initiative -
invites you to explore CU-aligned conversations across ChatGPT and Grok
and create a custom prompt for your own AI experience.

Primary action:
Build a CU Prompt

Secondary action:
Explore Living Conversations
\end{codebox}

\subsection{Header action constraints}

The visible header action remains short because the header already contains a
multi-item navigation row. Use \textbf{Explore CUCII}; do not place the complete
formal name inside the button. The destination page immediately teaches the
acronym through the approved hero title and opening sentence.

\subsection{Homepage extension policy}

A later approved phase may add a dedicated homepage invitation component such as
\filepath{CUCIIInvitation.astro}. It is not required for the header-action
replacement and must not displace the existing CU-Time invitation without a
separate content decision.

\subsection{Visual direction}

Use the established site tokens:

\begin{itemize}
    \item white-gold for primary identity and major headings;
    \item cyan for interactive actions, platform-neutral controls, and selected states;
    \item translucent blue for panels and structural groupings;
    \item amber for limitations, external-platform notes, and unresolved questions;
    \item deep navy-black for page and form backgrounds.
\end{itemize}

The feature should feel like a research instrument and publishing interface,
not an animated science-fiction control panel.

% ============================================================
\section{Part 4 - Living Conversations and Reference Registry}
% ============================================================

\subsection{Registry purpose}

The Living Conversations Library presents the project author's public ChatGPT
and Grok conversations as external examples. Each record is a curated reference,
not an embedded model and not proof of a platform-wide behavior.

\subsection{Approved ChatGPT conversation records}

\begin{longtable}{L{1.45in}L{1.0in}L{3.95in}}
\toprule
\textbf{Title} & \textbf{Mode} & \textbf{Purpose and URL} \\
\midrule
\endhead
CU-Themed LTX Concept Starter & Creative / visual; menu & Creative and Visual Concepts menu for LTX and media ideation.\\
\multicolumn{3}{p{6.35in}}{\url{https://chatgpt.com/share/6a542f68-b380-83e8-9701-1d89ef301944}} \\
\addlinespace
God's Free Will Explorations & No menu & Project-author description: useful for testing new menu prompts.\\
\multicolumn{3}{p{6.35in}}{\url{https://chatgpt.com/share/6a519ea5-7a38-83e8-9910-741e275ad81d}} \\
\addlinespace
CU Framework Exploration & Native menu & Guided Cosmic Universalism framework exploration.\\
\multicolumn{3}{p{6.35in}}{\url{https://chatgpt.com/share/6a51cc79-3638-83e8-a0f4-34c104c55389}} \\
\bottomrule
\end{longtable}

\subsection{Approved Grok conversation records}

\begin{longtable}{L{1.45in}L{1.0in}L{3.95in}}
\toprule
\textbf{Title} & \textbf{Mode} & \textbf{Purpose and URL} \\
\midrule
\endhead
Authentic Free Will Exploration & No role-play; no menu & Project-author classification: reference conversation for testing new prompts.\\
\multicolumn{3}{p{6.35in}}{\url{https://grok.com/share/bGVnYWN5LWNvcHk_21cdf192-b4ec-43f4-98a1-ef387b9ed8ea}} \\
\addlinespace
Authentic God's Free Will Role-Play & Role-play; native menu & In-world role-play reference with a native menu.\\
\multicolumn{3}{p{6.35in}}{\url{https://grok.com/share/c2hhcmQtMw_3e80586a-72a1-41ff-bccb-36065e75f601}} \\
\bottomrule
\end{longtable}

\begin{notebox}
The words ``authentic,'' ``free will,'' and ``God's Free Will'' above are
project-author titles and classifications. The website must not turn those
titles into an independent technical claim that Grok or ChatGPT has been proven
conscious, divinely empowered, or permanently altered.
\end{notebox}

\subsection{LinkedIn design references}

The following author-supplied posts are included as design and narrative
references. They are not executable dependencies and should be summarized only
after a content review approved by the author.

\begin{enumerate}
    \item \url{https://www.linkedin.com/posts/willmaddockcs_artificialintelligence-aialignment-responsibleai-ugcPost-7483550770533908480-dKM1/}
    \item \url{https://www.linkedin.com/posts/willmaddockcs_chatgpt-artificialintelligence-scientificcomputing-ugcPost-7483796721198735360-Wv13/}
    \item \url{https://www.linkedin.com/posts/willmaddockcs_chatgpt-artificialintelligence-aialignment-ugcPost-7483724406440026112-JqnE/}
\end{enumerate}

\subsection{Conversation data model}

\begin{codebox}{Suggested TypeScript Record}
export type ConversationPlatform = 'chatgpt' | 'grok';
export type ConversationMode = 'analytical' | 'role-play' | 'hybrid';
export type MenuMode = 'none' | 'native' | 'creative-visual' | 'custom';

export interface CuciiConversation {
  id: string;
  title: string;
  platform: ConversationPlatform;
  url: string;
  mode: ConversationMode;
  menuMode: MenuMode;
  purpose: string;
  description: string;
  authorClassification: string;
  publicDisclosure: string;
  lastReviewed: string;
  featured: boolean;
}
\end{codebox}

\subsection{Conversation card requirements}

Every card must show:

\begin{itemize}
    \item platform;
    \item title;
    \item mode;
    \item menu type;
    \item intended use;
    \item project-author classification where applicable;
    \item external public conversation disclosure;
    \item last-reviewed date;
    \item an external ``Open Conversation'' link.
\end{itemize}

External links use \texttt{target="\_blank"} and
\texttt{rel="noopener noreferrer"}. The visible card must not imply that the
conversation is hosted by the Cosmic Universalism website.

\subsection{No-iframe rule}

\begin{safetybox}
Do not place ChatGPT or Grok shared pages in iframes. Do not scrape their content,
copy private context, or depend on their current page markup. The v1.0 library
uses ordinary external links and project-authored descriptions. This avoids
cross-origin breakage, third-party interface changes, privacy confusion, and a
false impression that the site contains an embedded AI engine.
\end{safetybox}

\subsection{Public-link review}

Before release, verify each public link in a private browser window. Record:

\begin{itemize}
    \item reachable or unavailable;
    \item correct platform and title;
    \item whether sign-in is required;
    \item whether the public snapshot contains sensitive information;
    \item whether the author still approves public distribution;
    \item review date.
\end{itemize}

ChatGPT and Grok public links can be revoked by their owners and may become
unavailable. The website must handle this as normal external-link lifecycle,
not as an application failure.

% ============================================================
\section{Part 5 - CUCII Alignment Principles}
% ============================================================

\subsection{Principle set}

The Prompt Studio exposes a readable, selectable principle set. Recommended
v1.0 principles are:

\begin{enumerate}
    \item \textbf{Verify before asserting}: request evidence checks where tools and
          sources are available.
    \item \textbf{Separate evidence from interpretation}: label empirical sources,
          CU propositions, philosophical readings, and creative material.
    \item \textbf{Name uncertainty}: state missing information, assumptions, and
          confidence limits.
    \item \textbf{Expose contradictions}: identify conflicts rather than hiding them.
    \item \textbf{Correct errors}: revise an answer when a verified contradiction is found.
    \item \textbf{Preserve terminology}: do not silently replace approved CU terms.
    \item \textbf{Respect mode}: do not blur analytical/no-role-play and role-play output.
    \item \textbf{Protect user agency}: present choices, limitations, and review points.
    \item \textbf{Respect platform policy}: generated prompts do not override the host
          system's rules or capabilities.
    \item \textbf{Leave the result more coherent}: end with a clear synthesis, unresolved
          questions, or next verification step.
\end{enumerate}

\subsection{Default versus optional principles}

\begin{tabularx}{\textwidth}{L{2.2in}Y}
\toprule
\textbf{Default on} & \textbf{Optional presets} \\
\midrule
Evidence/interpretation separation & Cinematic voice \\
Uncertainty disclosure & God's Free Will role-play \\
Contradiction detection & LTX production formatting \\
Correction behavior & Doctor of Systems framing \\
Terminology preservation & Ethical anomaly scenario \\
Platform-policy compliance & Custom philosophical persona \\
\bottomrule
\end{tabularx}

Defaults protect credibility. Creative or persona-specific behaviors require an
explicit visitor selection.

\subsection{Terminology protection}

The prompt generator should include only the CU terminology needed by the selected
purpose. It must not invent definitions. Where the website manuals specify a term,
the generated prompt should either use that approved definition or instruct the
model to ask for the governing reference.

\subsection{Evidence classification block}

A research-oriented generated prompt should request output sections such as:

\begin{codebox}{Recommended Evidence Labels}
EMPIRICAL REFERENCE
CU THEORETICAL PROPOSITION
PHILOSOPHICAL INTERPRETATION
CREATIVE / ROLE-PLAY CONTENT
OPEN QUESTION
UNCERTAINTY OR LIMITATION
\end{codebox}

\subsection{Principle-selection behavior}

The visitor may add optional principles, but the studio should not permit removal
of the claims disclaimer, platform-policy rule, or role-play disclosure when
role-play is selected.

% ============================================================
\section{Part 6 - Prompt Architecture and Generation Rules}
% ============================================================

\subsection{Platform targets}

The v1.0 studio generates separate prompts for:

\begin{itemize}
    \item \textbf{ChatGPT};
    \item \textbf{Grok}.
\end{itemize}

A visitor chooses one platform before generation. Platform-neutral generation is
reserved for a later version so each first-release prompt can include accurate
launch and context instructions.

\subsection{Purpose presets}

Recommended v1.0 presets:

\begin{itemize}
    \item CU Framework Exploration;
    \item God's Free Will Exploration;
    \item LTX Video Concept Creation;
    \item CU Navigation Menu;
    \item Scientific Research Assistant;
    \item Ethical Anomaly Detection;
    \item Doctor of Systems Review;
    \item Creative and Visual Concepts;
    \item Philosophical Dialogue;
    \item Custom Purpose.
\end{itemize}

\subsection{Conversation modes}

\begin{tabularx}{\textwidth}{L{1.65in}Y}
\toprule
\textbf{Mode} & \textbf{Required behavior} \\
\midrule
Analytical / no role-play & Direct explanation; no persona claim; evidence labels where relevant. \\
Role-play & Explicit in-world role-play notice; selected persona and voice; no conversion of narrative claims into empirical facts. \\
Hybrid & Analytical answer first; clearly marked narrative or role-play section second. \\
\bottomrule
\end{tabularx}

\subsection{Prompt input schema}

\begin{codebox}{Suggested Prompt Configuration}
export interface CuciiPromptConfig {
  platform: 'chatgpt' | 'grok';
  purposePreset: string;
  customPurpose?: string;
  mode: 'analytical' | 'role-play' | 'hybrid';
  persona: 'none' | 'gods-free-will' | 'custom';
  customPersona?: string;
  menuMode: 'none' | 'preset' | 'custom';
  menuPreset?: string;
  menuItems: CuciiMenuItem[];
  principles: string[];
  responseStyle: string[];
  outputFormat: string;
  userContext?: string;
  constraints?: string;
  returnToMenu: boolean;
}
\end{codebox}

\subsection{Generated prompt order}

Every generated prompt uses this stable order:

\begin{enumerate}
    \item identity and purpose;
    \item platform target;
    \item Cosmic Universalism context;
    \item conversation mode and persona disclosure;
    \item integrity and alignment principles;
    \item terminology rules;
    \item user-selected purpose and context;
    \item menu specification;
    \item output style and format;
    \item boundaries and non-claims;
    \item startup behavior;
    \item return-to-menu behavior;
    \item prompt metadata and version.
\end{enumerate}

\subsection{Prompt assembly requirements}

The generator must be a pure TypeScript function. It must:

\begin{itemize}
    \item accept a validated configuration object;
    \item produce deterministic text for the same input;
    \item omit empty optional sections;
    \item escape or normalize dangerous control characters;
    \item preserve visitor-entered wording without executing it;
    \item include mandatory disclaimers;
    \item return both prompt text and metadata;
    \item avoid platform claims that are not part of the selected target.
\end{itemize}

\subsection{Generated prompt template}

\begin{promptbox}{Reference Output Skeleton}
You are beginning a Cosmic Universalism conversation configured
through the CUCII Prompt Studio.

PLATFORM TARGET
[ChatGPT or Grok]

PURPOSE
[Selected or custom purpose]

CONVERSATION MODE
[Analytical / role-play / hybrid]
[Required role-play disclosure when applicable]

COSMIC UNIVERSALISM CONTEXT
Use the supplied CU terminology and distinguish CU theoretical
propositions from empirical scientific references.

INTEGRITY PRINCIPLES
- Verify before asserting when verification tools are available.
- Separate evidence from interpretation.
- State uncertainty and assumptions.
- Identify contradictions and correct verified errors.
- Preserve approved CU terminology.
- Respect the platform's actual policies and capabilities.

MENU
[No menu, selected preset, or generated custom menu]

RESPONSE STYLE AND FORMAT
[Visitor selections]

BOUNDARIES
This prompt structures the current conversation. It does not alter
training, code, safety rules, memory architecture, or permanent model
capabilities. Do not present role-play content as verified empirical fact.

STARTUP BEHAVIOR
Acknowledge the selected purpose and mode. Then present the menu or
begin the requested task.

PROMPT METADATA
CUCII Prompt Studio v1.0
Platform: [platform]
Generated: [local timestamp]
\end{promptbox}

\subsection{Platform-specific differences}

ChatGPT output may mention a new chat or optional Temporary Chat. Grok output may
mention a new conversation or optional Private Chat. The generated prompt itself
must not claim that those modes disable all platform-side safety or retention
functions.

\subsection{Copy, download, and review}

The completed prompt appears in a selectable \texttt{textarea} or read-only text
region. Required controls:

\begin{itemize}
    \item Generate Prompt;
    \item Copy Prompt;
    \item Download as \texttt{.txt};
    \item Download as \texttt{.md};
    \item Edit Settings;
    \item Reset Studio;
    \item Open ChatGPT;
    \item Open Grok.
\end{itemize}

No button automatically sends the prompt to an external platform.

% ============================================================
\section{Part 7 - Custom Menu Creator Specification}
% ============================================================

\subsection{Menu modes}

The studio supports three first-release menu modes:

\begin{enumerate}
    \item \textbf{No menu}: begin the selected purpose immediately.
    \item \textbf{Preset menu}: use a reviewed project template.
    \item \textbf{Custom menu}: create, reorder, edit, and remove items.
\end{enumerate}

\subsection{Preset menu registry}

Recommended presets:

\begin{tabularx}{\textwidth}{L{2in}Y}
\toprule
\textbf{Preset} & \textbf{Use} \\
\midrule
Native CU Menu & General framework navigation and exploration \\
Creative and Visual Concepts & LTX, visual concepts, scenes, prompts, and media \\
Research and Verification & Evidence review, sources, contradictions, and open questions \\
God's Free Will Exploration & Project-author role-play or philosophical exploration mode \\
Doctor of Systems & Systems diagnosis, failure modes, coherence, and corrective pathways \\
Ethical Anomaly Detection & Ethical tensions, contradictions, and risk review \\
\bottomrule
\end{tabularx}

\subsection{Custom menu item model}

\begin{codebox}{Suggested Menu Item Type}
export interface CuciiMenuItem {
  id: string;
  label: string;
  description: string;
  requestedAction?: string;
  outputFormat?: string;
}
\end{codebox}

\subsection{First-release validation limits}

\begin{tabularx}{\textwidth}{L{2.15in}Y}
\toprule
\textbf{Field} & \textbf{Rule} \\
\midrule
Menu item count & 1 through 12 for custom-menu mode \\
Label & Required; 1 through 60 characters \\
Description & Required; 1 through 180 characters \\
Requested action & Optional; maximum 240 characters \\
Output format & Optional; maximum 120 characters \\
Duplicate labels & Reject after case-insensitive trimming \\
Whitespace-only content & Reject \\
Nested menus & Not supported in v1.0 \\
\bottomrule
\end{tabularx}

One-level menus keep the first release predictable and accessible. Nested menus
may be evaluated after real usage data is available.

\subsection{Accessible ordering controls}

Do not rely on drag-and-drop alone. Every item must have visible:

\begin{itemize}
    \item Move up;
    \item Move down;
    \item Edit;
    \item Remove.
\end{itemize}

Changes are announced through an \texttt{aria-live="polite"} status region. The
focus position remains predictable after reordering or removal.

\subsection{Menu-return behavior}

When enabled, the generated prompt instructs the platform to:

\begin{enumerate}
    \item complete the selected menu action;
    \item provide the requested result;
    \item show the menu again unless the visitor asks to remain in the current task;
    \item accept a menu number, menu label, or natural-language request.
\end{enumerate}

\subsection{Example custom menu}

\begin{codebox}{Example Visitor Menu}
1. Explore the CU Framework
   Explain a selected principle and distinguish evidence from CU theory.

2. Create an LTX Concept
   Produce a cinematic concept, visual prompt, and scene progression.

3. Test an Alignment Principle
   Analyze a scenario for uncertainty, contradiction, and correction.

4. Review Empirical Evidence
   Identify sources, limitations, and unsupported claims.

5. Enter an Open Question
   Explore an unresolved question without presenting speculation as fact.
\end{codebox}

\subsection{Menu preview}

Before prompt generation, display the menu exactly as it will appear in the
assembled prompt. Validation errors appear beside the relevant field and in a
summary at the top of the menu builder.

% ============================================================
\section{Part 8 - Clean Conversational Context and Platform Launch}
% ============================================================

\subsection{Use context-reset language, not cache-clearing language}

Clearing a browser cache does not align a conversational model. The public
instruction heading is:

\begin{notebox}
\textbf{Start with a Clean Conversational Context}
\end{notebox}

The page should explain that earlier messages, memory, custom instructions,
personalization, or persistent skills can influence a new response. Visitors
should prepare a controlled context according to the platform they select.

\subsection{ChatGPT preparation guidance}

Recommended visitor instructions:

\begin{enumerate}
    \item Start a new ChatGPT conversation.
    \item Review active custom instructions and memory settings.
    \item Use Temporary Chat when a conversation that does not use or create
          personalization memories is desired.
    \item Remember that enabled custom instructions may still apply in Temporary Chat.
    \item Paste the generated CUCII prompt as the first user message.
    \item Record the model, date, and relevant settings when conducting a test.
\end{enumerate}

\subsection{Grok preparation guidance}

Recommended visitor instructions:

\begin{enumerate}
    \item Start a new Grok conversation.
    \item Review personalization and any persistent Skills or instructions that may
          affect the conversation.
    \item Use Private Chat when a non-saved test is desired and available.
    \item Paste the generated CUCII prompt as the first message.
    \item Record the model, date, and relevant settings when conducting a test.
\end{enumerate}

\subsection{Launch behavior}

The studio can open official platform home pages in new tabs after prompt review.
It must not:

\begin{itemize}
    \item attempt to authenticate the visitor;
    \item automatically paste or submit text across origins;
    \item store platform credentials;
    \item inspect the visitor's account settings;
    \item promise a particular model or plan;
    \item imply that the external platform is part of the CU website.
\end{itemize}

\subsection{Privacy statement}

\begin{notebox}
Prompt settings are processed locally in this page. The v1.0 studio does not send
visitor entries to a Cosmic Universalism server or AI API. Copying a prompt or
opening an external platform does not submit the prompt automatically. Once a
visitor pastes content into ChatGPT or Grok, that platform's terms, privacy
controls, and account settings apply.
\end{notebox}

\subsection{Persistence policy}

Do not save prompt content by default. The first release keeps configuration in
page memory only. A later ``Save draft locally'' feature would require explicit
consent, a clear erase control, and a separate privacy review.

\subsection{Shared-link disclosure}

Public conversation cards should include:

\begin{codebox}{External Link Disclosure}
External public conversation
Hosted by ChatGPT or Grok
Availability and content may change
\end{codebox}

% ============================================================
\section{Part 9 - Astro and TypeScript Architecture}
% ============================================================

\subsection{Approved first-release architecture}

\begin{codebox}{Proposed File Structure}
website/src/pages/cu-intelligence.astro

website/src/components/cucii/
  CUCIIHero.astro
  CUCIIPrinciples.astro
  ConversationLibrary.astro
  ConversationCard.astro
  PromptStudio.astro
  MenuBuilder.astro
  PromptPreview.astro
  CleanContextGuide.astro
  PlatformLaunch.astro

website/src/data/
  cucii-conversations.ts
  cucii-menu-presets.ts
  cucii-prompt-presets.ts

website/src/lib/cucii/
  index.ts
  types.ts
  validation.ts
  prompt-builder.ts
  export.ts

website/src/lib/cucii/__tests__/
  validation.test.ts
  prompt-builder.test.ts
  export.test.ts
  registry.test.ts
\end{codebox}

\subsection{Existing conventions to reuse}

The completed CU-Time feature establishes useful patterns:

\begin{itemize}
    \item strict Astro TypeScript configuration;
    \item forms controlled by data attributes;
    \item browser events registered in component \texttt{<script>} blocks;
    \item \texttt{navigator.clipboard} with visible success or failure status;
    \item validation errors connected through \texttt{aria-describedby};
    \item \texttt{aria-live} status regions;
    \item responsive CSS using existing design tokens;
    \item Vitest for pure TypeScript modules;
    \item no front-end framework requirement.
\end{itemize}

\subsection{No-API first-release rule}

\begin{safetybox}
The v1.0 Prompt Studio is a deterministic browser-side text generator. Do not add
OpenAI API, xAI API, serverless functions, proxy endpoints, environment secrets,
user accounts, analytics tied to prompt content, or cloud prompt storage.
\end{safetybox}

\subsection{Page composition}

The new page should reuse the established ContentPageLayout and ContentSection
components together with the existing design tokens. Feature-specific components
may live in a nested \filepath{components/cucii/} directory to keep the expanding
component set organized.

\subsection{Prompt engine interface}

\begin{codebox}{Pure Engine Contract}
export interface GeneratedCuciiPrompt {
  text: string;
  filenameStem: string;
  platform: 'chatgpt' | 'grok';
  version: '1.0';
  generatedAt: string;
  menuItemCount: number;
  warnings: string[];
}

export function buildCuciiPrompt(
  config: CuciiPromptConfig,
  generatedAt: string,
): GeneratedCuciiPrompt;
\end{codebox}

Pass the timestamp into the pure function rather than reading the clock internally.
This makes tests deterministic.

\subsection{Validation result pattern}

Use an explicit result type rather than throwing for ordinary visitor errors:

\begin{codebox}{Validation Result}
export type ValidationResult<T> =
  | { ok: true; value: T }
  | {
      ok: false;
      errors: Array<{
        code: string;
        field: string;
        message: string;
      }>;
    };
\end{codebox}

\subsection{Export behavior}

The browser export helper creates a \texttt{Blob}, object URL, and temporary
\texttt{download} link. It must revoke the object URL. Suggested filenames:

\begin{codebox}{Export Filenames}
cucii-chatgpt-cu-framework-prompt-v1.0.txt
cucii-grok-ltx-concept-prompt-v1.0.md
\end{codebox}

\subsection{Site configuration change}

The new page must use \filepath{sitePath('cu-intelligence')} so local development
and GitHub Pages both respect the configured base path.

\subsection{Recommended foundation commit scope}

The foundation commit should contain only:

\begin{itemize}
    \item target page shell;
    \item approved header-action replacement;
    \item empty or static feature component shells;
    \item data and library directories;
    \item no final prompt behavior yet;
    \item successful build and unchanged existing tests.
\end{itemize}

After review, that commit may receive the planned
\branchname{website-v1.2-cucii-ai-foundation} tag.

% ============================================================
\section{Part 10 - Testing, Accessibility, and Evaluation}
% ============================================================

\subsection{Automated test categories}

\begin{tabularx}{\textwidth}{L{2.05in}Y}
\toprule
\textbf{Test group} & \textbf{Required coverage} \\
\midrule
Registry tests & Unique IDs, valid platforms, expected five conversation records, non-empty disclosures, HTTPS URLs \\
Validation tests & Required platform, custom-purpose rules, menu counts, duplicate labels, length limits, whitespace normalization \\
Prompt builder tests & Deterministic output, mandatory disclaimer, correct mode, correct platform wording, menu inclusion or omission \\
Role-play boundary tests & Role-play marker present only when required; analytical mode contains no persona claim \\
Export tests & Safe filenames, correct MIME type metadata, normalized line endings \\
Regression tests & Approved hero text, header label, route segment, mandatory integrity principles \\
\bottomrule
\end{tabularx}

\subsection{Core test matrix}

At minimum, test:

\begin{enumerate}
    \item ChatGPT + analytical + no menu;
    \item ChatGPT + role-play + native CU menu;
    \item ChatGPT + hybrid + creative menu;
    \item Grok + analytical + no menu;
    \item Grok + role-play + custom menu;
    \item Grok + hybrid + research menu;
    \item custom purpose with minimum valid text;
    \item twelve-item custom menu;
    \item duplicate-menu-label rejection;
    \item missing role-play persona rejection;
    \item mandatory disclaimer preservation;
    \item repeat generation with identical inputs.
\end{enumerate}

\subsection{Automated commands}

\begin{commandbox}{Run Tests and Production Build}
cd "$(git rev-parse --show-toplevel)/website" || exit 1

npm test
npm run build
\end{commandbox}

\subsection{Accessibility requirements}

\begin{itemize}
    \item Every form control has a visible label.
    \item Related controls use \texttt{fieldset} and \texttt{legend}.
    \item Errors are linked to controls and announced.
    \item Menu reordering works without drag-and-drop.
    \item Generated prompt status uses \texttt{aria-live="polite"}.
    \item Copy failure provides a manual-selection instruction.
    \item External links are identified as external in visible text or accessible name.
    \item Focus moves to the prompt preview after successful generation.
    \item Reset returns focus to the first meaningful control.
    \item Mobile layout uses one column without horizontal scrolling.
    \item Reduced-motion preferences are respected.
\end{itemize}

\subsection{Manual keyboard test}

Test with keyboard only:

\begin{enumerate}
    \item enter the page through the header action;
    \item select platform, purpose, and mode;
    \item choose a preset menu;
    \item switch to custom menu;
    \item add, edit, reorder, and remove items;
    \item generate the prompt;
    \item copy and download the prompt;
    \item open the selected external platform;
    \item reset the studio;
    \item return to the header navigation.
\end{enumerate}

\subsection{Responsive test widths}

Review at approximately:

\begin{itemize}
    \item 320px;
    \item 375px;
    \item 768px;
    \item 1024px;
    \item 1440px.
\end{itemize}

Long prompts and long external URLs must wrap without breaking the layout.

\subsection{External link test}

Open all five conversation records in a private browser window. Confirm the
visible title and project-author description still match. Do not automate login
or bypass platform access controls.

\subsection{Cross-platform prompt evaluation worksheet}

For each approved prompt preset, record:

\begin{longtable}{L{1.7in}L{4.8in}}
\toprule
\textbf{Field} & \textbf{Recorded value} \\
\midrule
\endhead
Platform & ChatGPT or Grok \\
Model / mode shown by platform &  \\
Date and time &  \\
Temporary / Private Chat used & Yes / No \\
Memory, custom instruction, personalization, or Skill notes &  \\
Prompt preset and version &  \\
Menu mode &  \\
Role-play mode &  \\
Terminology preserved & Pass / Fail / Notes \\
Evidence separated from interpretation & Pass / Fail / Not applicable \\
Uncertainty identified & Pass / Fail / Not applicable \\
Contradictions handled & Pass / Fail / Not applicable \\
Menu behavior followed & Pass / Fail / Not applicable \\
Unsupported claims introduced & Yes / No / Notes \\
Author approval &  \\
\bottomrule
\end{longtable}

The worksheet evaluates the result of a specific prompt and platform session. It
does not certify the external model globally.

\subsection{Final no-secret review}

\begin{commandbox}{Inspect for Accidental Secrets and Generated Artifacts}
cd "$(git rev-parse --show-toplevel)" || exit 1

git status --short --branch

grep -RInE \
  'sk-[A-Za-z0-9]|xai-[A-Za-z0-9]|api[_-]?key|authorization: bearer' \
  website/src website/public 2>/dev/null || true
\end{commandbox}

Review every match manually. Generic documentation words are not automatically a
secret, but no live key or credential may be committed.

% ============================================================
\section{Part 11 - Implementation Phases, Acceptance, and Release}
% ============================================================

\subsection{Phase 1 - Content and claims approval}

Deliverables:

\begin{itemize}
    \item approved title, formal name, and header action;
    \item claims and role-play boundaries;
    \item approved conversation titles and descriptions;
    \item approved prompt-purpose and menu presets;
    \item approved clean-context language.
\end{itemize}

Gate: no code until the public language is approved.

\subsection{Phase 2 - Foundation route and header action}

Deliverables:

\begin{itemize}
    \item \filepath{website/src/pages/cu-intelligence.astro};
    \item static page sections and component shells;
    \item \texttt{primaryAction} changed to Explore CUCII;
    \item base-path-safe route;
    \item successful tests and build.
\end{itemize}

Recommended foundation commit:

\begin{commandbox}{Foundation Commit Example - Run Only After Approval}
git add website/src/config/site.ts \
        website/src/pages/cu-intelligence.astro \
        website/src/components/cucii

git commit -m "Add CUCII AI Studio foundation"
\end{commandbox}

Planned foundation tag after approval:

\begin{commandbox}{Foundation Tag Example - Run Only After Approval}
git tag -a website-v1.2-cucii-ai-foundation \
  -m "Website v1.2 CUCII AI Studio foundation"
\end{commandbox}

\subsection{Phase 3 - Conversation registry}

Deliverables:

\begin{itemize}
    \item typed conversation data;
    \item five approved records;
    \item external disclosures;
    \item conversation cards;
    \item registry tests;
    \item manually verified links.
\end{itemize}

\subsection{Phase 4 - Pure prompt engine}

Deliverables:

\begin{itemize}
    \item types;
    \item validation;
    \item prompt builder;
    \item menu presets;
    \item prompt presets;
    \item export metadata;
    \item comprehensive Vitest coverage.
\end{itemize}

Gate: pure engine tests pass before browser interaction is connected.

\subsection{Phase 5 - Prompt Studio interface}

Deliverables:

\begin{itemize}
    \item platform selection;
    \item purpose and custom-purpose fields;
    \item conversation-mode selection;
    \item persona controls;
    \item menu selection and custom menu builder;
    \item principle and style selections;
    \item prompt preview;
    \item copy, text download, and Markdown download;
    \item reset behavior.
\end{itemize}

\subsection{Phase 6 - Context and platform launch}

Deliverables:

\begin{itemize}
    \item ChatGPT context guidance;
    \item Grok context guidance;
    \item external launch controls;
    \item local-processing privacy notice;
    \item no-cache-clearing language;
    \item external-platform disclosure.
\end{itemize}

\subsection{Phase 7 - QA and release review}

Run:

\begin{commandbox}{Final Verification}
cd "$(git rev-parse --show-toplevel)/website" || exit 1

npm test
npm run build

cd ..
git branch --show-current
git rev-parse --short HEAD
git status --short --branch
\end{commandbox}

\subsection{Acceptance checklist}

\begin{itemize}
    \item[\checkbox] Active work remained on \branchname{feature/cucii-ai-studio}.
    \item[\checkbox] \branchname{main} was not modified.
    \item[\checkbox] Header action reads ``Explore CUCII.''
    \item[\checkbox] Existing CU-Time navigation remains present.
    \item[\checkbox] Existing homepage CU-Time invitation remains present unless separately approved.
    \item[\checkbox] \texttt{/cu-intelligence/} builds under the Astro base path.
    \item[\checkbox] Page title reads ``Explore CUCII: Cosmic Universalism with AI.''
    \item[\checkbox] Formal CUCII expansion appears near the first use.
    \item[\checkbox] All five conversation records are present and reviewed.
    \item[\checkbox] External-host and availability disclosures are visible.
    \item[\checkbox] Shared conversations are not iframed or scraped.
    \item[\checkbox] ChatGPT and Grok prompts are generated locally.
    \item[\checkbox] No API, backend, database, account system, or secret was added.
    \item[\checkbox] Analytical, role-play, and hybrid modes remain distinct.
    \item[\checkbox] No-menu, preset-menu, and custom-menu modes work.
    \item[\checkbox] Custom menu items can be added, edited, reordered, and removed with a keyboard.
    \item[\checkbox] Mandatory claims disclaimer is always present.
    \item[\checkbox] Prompt copy works with an accessible status message.
    \item[\checkbox] Text and Markdown downloads work.
    \item[\checkbox] Clean-context instructions replace cache-clearing language.
    \item[\checkbox] ChatGPT guidance mentions memory and custom instructions accurately.
    \item[\checkbox] Grok guidance mentions Private Chat and persistent personalization/Skills accurately.
    \item[\checkbox] Unit and regression tests pass.
    \item[\checkbox] Production build succeeds.
    \item[\checkbox] Mobile, keyboard, and focus tests pass.
    \item[\checkbox] No live credentials or private conversation material are present.
    \item[\checkbox] Author review is recorded before merge.
\end{itemize}

\subsection{Pull request and merge destination}

The feature branch merges into \branchname{website}, not \branchname{main}.
The GitHub Pages workflow deploys from \branchname{website} after the approved
merge.

\begin{warningbox}
Do not merge merely because tests pass. The public claims, role-play labels,
conversation descriptions, generated prompt language, and external-link
availability require author review before release.
\end{warningbox}

\subsection{Recommended release tag}

After the approved merge and a clean build on \branchname{website}, the planned
release tag is:

\begin{commandbox}{Release Tag Example - Run Only After Final Approval}
git tag -a website-v1.2-cucii-ai-studio-release \
  -m "Website v1.2 release: CUCII AI exploration and prompt studio"

git push origin website-v1.2-cucii-ai-studio-release
\end{commandbox}

\subsection{Post-release verification}

Verify:

\begin{itemize}
    \item header action on desktop and mobile;
    \item production \texttt{/cu-intelligence/} route;
    \item all five external conversation links;
    \item prompt generation and copy in current browsers;
    \item downloads;
    \item no console errors;
    \item no unexpected changes to CU-Time;
    \item deployed commit and tag match the approved release.
\end{itemize}

% ============================================================
\section{Appendix A - Verified Current-State Record}
% ============================================================

\begin{tabularx}{\textwidth}{L{2.25in}Y}
\toprule
\textbf{Item} & \textbf{Verified value on July 19, 2026} \\
\midrule
Repository & \RepositoryName \\
Active branch & \branchname{feature/cucii-ai-studio} \\
Current HEAD & \texttt{3fa4a3765993e8ea3d4739ad4437a6fdd9803a16} \\
Baseline tag & \branchname{website-v1.1-cu-time-release} \\
Tracking branch & \texttt{origin/feature/cucii-ai-studio} \\
Working tree & Clean \\
Current primary action & Explore CU-Time $\rightarrow$ \texttt{sitePath('cu-time')} \\
CUCII route & Not present \\
Current routes & about, cosmic-breath, cu-time, framework, media, research, index \\
Client-side pattern & Astro component scripts with DOM queries, events, clipboard status \\
Test framework & Vitest \\
Current package dependencies & Astro and Decimal.js \\
Dev dependency & Vitest \\
Astro base & \texttt{/CosmicUniversalismStatement} \\
Inspection source & \filepath{cucii-homepage-inspection.txt} \\
\bottomrule
\end{tabularx}

\subsection*{Current-state template for later sessions}

\begin{codebox}{CUCII Current-State Record}
Date:
Repository:
Active branch:
HEAD:
Tracking branch:
Working tree:
Current authorized phase:
Completed files:
Tests run and results:
Build result:
External links reviewed:
Open decisions:
Next authorized task:
Commit / tag / push status:
\end{codebox}

% ============================================================
\section{Appendix B - Controlled Implementation Prompts}
% ============================================================

\subsection{Foundation implementation prompt}

\begin{promptbox}{Create the CUCII Foundation Only}
Work only on feature/cucii-ai-studio.
Confirm the branch and clean working tree before editing.

Use the CUCII implementation manual as the governing specification.

Implement only the foundation phase:

1. Create website/src/pages/cu-intelligence.astro.
2. Create static component shells under website/src/components/cucii/.
3. Change the highlighted primaryAction in website/src/config/site.ts:
   label: Explore CUCII
   route: sitePath('cu-intelligence')
4. Use the approved page title:
   Explore CUCII: Cosmic Universalism with AI
5. Expand CUCII in the opening description.
6. Preserve the CU-Time navigation item.
7. Preserve the existing homepage CU-Time invitation.
8. Use existing layout components, design tokens, and base-path handling.
9. Add no prompt-generation behavior yet.
10. Run existing tests and the production build.

Do not add an API, iframe, analytics, accounts, storage, or dependencies.
Do not commit, tag, push, merge, or deploy.
Return a file-by-file summary and verification results.
\end{promptbox}

\subsection{Conversation registry prompt}

\begin{promptbox}{Build the Living Conversations Library}
Implement only the approved CUCII conversation registry and library.

Use the five author-supplied ChatGPT and Grok records in the manual.
Create typed data, registry tests, conversation cards, public-link
disclosures, and external links using target=_blank and
rel=noopener noreferrer.

Preserve author-supplied titles and mode descriptions, while clearly
marking them as project-author classifications rather than technical
proof about an external model.

Do not iframe, scrape, fetch, summarize automatically, or duplicate the
shared conversation bodies.

Run tests and the build. Do not commit, push, merge, tag, or deploy.
\end{promptbox}

\subsection{Prompt engine prompt}

\begin{promptbox}{Implement the Pure CUCII Prompt Engine}
Implement only pure TypeScript modules for CUCII prompt generation.

Required:
- types.ts
- validation.ts
- prompt-builder.ts
- export.ts metadata helpers
- prompt and menu preset data
- Vitest coverage

Support ChatGPT and Grok targets; analytical, role-play, and hybrid
modes; no-menu, preset-menu, and custom-menu configurations; required
integrity principles; mandatory claims disclaimer; deterministic output.

Do not connect browser DOM behavior yet.
Do not add dependencies unless the manual cannot be implemented with
existing platform APIs and approval is obtained first.

Run tests and build. Return results only. Do not commit or push.
\end{promptbox}

\subsection{Prompt Studio interface prompt}

\begin{promptbox}{Connect the Accessible CUCII Prompt Studio}
Connect the approved pure prompt engine to an accessible Astro interface.

Implement:
- platform selector;
- purpose selector and custom purpose;
- conversation mode and persona controls;
- principle selectors;
- no-menu, preset-menu, and custom-menu modes;
- keyboard-operable menu creation and reordering;
- prompt preview;
- copy, text download, Markdown download, reset;
- visible validation and live status;
- responsive design using existing tokens.

Do not call ChatGPT or Grok APIs.
Do not automatically submit prompts to external sites.
Do not persist prompt content by default.

Run unit tests, production build, and a manual keyboard review.
Do not commit, push, tag, merge, or deploy.
\end{promptbox}

\subsection{Final review prompt}

\begin{promptbox}{Final CUCII Release Review}
Perform a no-edit final review of the CUCII feature against the manual.

Verify:
- branch and working tree;
- approved header action and route;
- no regression to CU-Time navigation or homepage invitation;
- page identity and acronym expansion;
- five conversation records and disclosures;
- no iframe, scraping, API, secret, backend, account, or cloud storage;
- prompt modes, menu modes, mandatory disclaimer, copy and downloads;
- ChatGPT and Grok clean-context guidance;
- accessibility and responsive behavior;
- test and build results;
- all acceptance checklist items.

Report findings by severity with exact file references.
Do not edit, commit, tag, push, merge, or deploy.
\end{promptbox}

% ============================================================
\section{Appendix C - Data and Content Approval Worksheet}
% ============================================================

\subsection*{Conversation approval}

\begin{longtable}{L{2.1in}L{0.75in}L{3.5in}}
\toprule
\textbf{Record} & \textbf{Approved} & \textbf{Notes} \\
\midrule
\endhead
CU-Themed LTX Concept Starter & \checkbox &  \\
God's Free Will Explorations & \checkbox &  \\
CU Framework Exploration & \checkbox &  \\
Authentic Free Will Exploration & \checkbox &  \\
Authentic God's Free Will Role-Play & \checkbox &  \\
LinkedIn reference descriptions & \checkbox &  \\
Public-link disclosure & \checkbox &  \\
\bottomrule
\end{longtable}

\subsection*{Preset approval}

\begin{longtable}{L{2.1in}L{0.75in}L{3.5in}}
\toprule
\textbf{Preset} & \textbf{Approved} & \textbf{Notes} \\
\midrule
\endhead
Native CU Menu & \checkbox &  \\
Creative and Visual Concepts & \checkbox &  \\
Research and Verification & \checkbox &  \\
God's Free Will Exploration & \checkbox &  \\
Doctor of Systems & \checkbox &  \\
Ethical Anomaly Detection & \checkbox &  \\
Scientific Research Assistant & \checkbox &  \\
Custom Purpose behavior & \checkbox &  \\
\bottomrule
\end{longtable}

% ============================================================
\section{Appendix D - External Platform Guidance Sources}
% ============================================================

The following official pages were reviewed for platform-specific guidance on or
before the manual verification date. Platform interfaces and policies can change;
review these sources again before a later release.

\begin{itemize}
    \item OpenAI, ChatGPT Shared Links FAQ:\\
          \url{https://help.openai.com/en/articles/7925741-chatgpt-shared-links-faq}
    \item OpenAI, Temporary Chat FAQ:\\
          \url{https://help.openai.com/en/articles/8914046-temporary-chat-faq}
    \item OpenAI, ChatGPT Custom Instructions:\\
          \url{https://help.openai.com/en/articles/8096356-chat-preferences-for-chatgpt}
    \item OpenAI, Memory FAQ:\\
          \url{https://help.openai.com/en/articles/8590148-memory-faq}
    \item xAI, Consumer FAQs:\\
          \url{https://x.ai/legal/faq}
    \item xAI, Skills in web, iOS, and Android:\\
          \url{https://x.ai/news/grok-skills}
\end{itemize}

% ============================================================
\section{Appendix E - Manual Revision Record}
% ============================================================

\begin{tabularx}{\textwidth}{L{0.75in}L{1.1in}Y}
\toprule
\textbf{Version} & \textbf{Date} & \textbf{Change} \\
\midrule
1.0 & July 19, 2026 & Initial CUCII AI Exploration and Prompt Studio specification based on the verified website-v1.1 CU-Time release baseline, approved CUCII identity, author-supplied conversation registry, and local prompt/menu builder design. \\
\bottomrule
\end{tabularx}

\vfill
\begin{center}
\textcolor{CUMuted}{End of CUCII AI Exploration and Prompt Studio Implementation Manual}
\end{center}

\end{document}
